% GENERATED FILE. Edit canonical lessons, then run npm run generate:books.
% canonical-source-hash: fnv1a64:a86124b691d8e37b
% canonical-lessons: LA-C298-evenio, LA-C298-fio, LA-C298-orior, LA-C298-intervenio, LA-C298-posterus, LA-R298-review-291-298

\chapter{Quid ēvēnit}
\label{ch:la-quid-evenit}
\hlchaptermodality{298}

\begin{chapteropening}
\textbf{By the end of this chapter:} \emph{I can say what happened and what came of it: to turn out, to become or happen, to arise, to come along and the next day.}

Complete the last lesson of chapter 298: I can say what happened and what came of it: to turn out, to become or happen, to arise, to come along and the next day.
\end{chapteropening}

\section[ēveniō, ēvenīre]{ēveniō, ēvenīre — to happen, to turn out}

\label{lesson:LA-C298-evenio}



\begin{quote}
Before the new one: say the Latin for to complete, then the Latin for to break off.
\end{quote}

\subsection*{You'll want to know: ēveniō, ēvenīre}
\textbf{ēveniō, ēvenīre} — \textquotedblleft{}to happen, to turn out\textquotedblright{}.

\begin{quote}
Mihi rēs mīra herī ēvēnit.
\end{quote}

\begin{quote}
\emph{A strange thing happened to me yesterday.}
\end{quote}

\emph{ēveniō, ēvenīre, ēvēnī, ēventum}, often with the dative: \emph{mihi ēvēnit}. Its noun \emph{ēventus} is taught.

\begin{cousinweb}
\textbf{ēveniō} is built on \emph{veniō}, \textquotedblleft{}to come\textquotedblright{}, with \emph{ex-}, \textquotedblleft{}out\textquotedblright{} (EtymDB): to come out, to turn out. English \emph{event} comes from its participle.
\end{cousinweb}

\begin{grammarlens}[title={What you've built}]
One more word for telling, explaining and describing at B1.
\end{grammarlens}

\subsection*{Your turn}
\begin{itemize}
\raggedright
  \item \emph{Say it:} \emph{ēveniō, ēvenīre}
  \item \emph{Say it:} \emph{ēveniō, ēvenīre}, once more
  \item \emph{Say it:} a sentence of your own with \emph{ēveniō, ēvenīre}
  \item \emph{Recall:} say the Latin for to complete, then the Latin for to break off, then say \emph{ēveniō, ēvenīre} again
\end{itemize}

\begin{tcolorbox}[breakable,colback=teal!4,colframe=teal!35!black,title={Before you move on}]
What does ēveniō, ēvenīre mean? (\textquotedblleft{}To happen, to turn out\textquotedblright{}.) Say it once more.
\end{tcolorbox}
\section[fīō, fierī]{fīō, fierī — to become; to happen, to be done}

\label{lesson:LA-C298-fio}



\begin{quote}
Before the new one: say the Latin for to break off, then the Latin for to happen.
\end{quote}

\subsection*{You'll want to know: fīō, fierī}
\textbf{fīō, fierī} — \textquotedblleft{}to become; to happen, to be done\textquotedblright{}.

\begin{quote}
Quid hīc fit? Cūr omnēs currunt?
\end{quote}

\begin{quote}
\emph{What is going on here? Why is everyone running?}
\end{quote}

\emph{fīō, fierī, factus sum}. It serves as the passive of \emph{faciō}: \emph{fit} 'it is done, it happens'.

\begin{cousinweb}
\textbf{fīō} is an inherited verb that serves as the passive of \emph{faciō} (Lewis \& Short). It is usually traced to the Proto-Indo-European root \emph{*b\textsuperscript{h}uH-}, \textquotedblleft{}to become, to be\textquotedblright{}, which also gives the Latin perfect \emph{fuī} and English \emph{be}.
\end{cousinweb}

\begin{grammarlens}[title={What you've built}]
One more word for telling, explaining and describing at B1.
\end{grammarlens}

\subsection*{Your turn}
\begin{itemize}
\raggedright
  \item \emph{Say it:} \emph{fīō, fierī}
  \item \emph{Say it:} \emph{fīō, fierī}, once more
  \item \emph{Say it:} a sentence of your own with \emph{fīō, fierī}
  \item \emph{Recall:} say the Latin for to break off, then the Latin for to happen, then say \emph{fīō, fierī} again
\end{itemize}

\begin{tcolorbox}[breakable,colback=teal!4,colframe=teal!35!black,title={Before you move on}]
What does fīō, fierī mean? (\textquotedblleft{}To become; to happen, to be done\textquotedblright{}.) Say it once more.
\end{tcolorbox}
\section[orior, orīrī]{orior, orīrī — to rise; to arise, to begin}

\label{lesson:LA-C298-orior}



\begin{quote}
Before the new one: say the Latin for to happen, then the Latin for to become.
\end{quote}

\subsection*{You'll want to know: orior, orīrī}
\textbf{orior, orīrī} — \textquotedblleft{}to rise; to arise, to begin\textquotedblright{}.

\begin{quote}
Sōle oriente ē portā discessimus.
\end{quote}

\begin{quote}
\emph{As the sun was rising, we left by the gate.}
\end{quote}

\emph{orior, orīrī, ortus sum}, a deponent. The sun rises, and so do quarrels and wars. The noun \emph{ortus}, 'sunrise', is taught. English \emph{orient}.

\begin{cousinweb}
\textbf{orior} is an inherited verb (EtymDB). English \emph{orient}, the land of the rising sun, comes from its participle, and \emph{origin} from its noun \emph{orīgō}.
\end{cousinweb}

\begin{grammarlens}[title={What you've built}]
One more word for telling, explaining and describing at B1.
\end{grammarlens}

\subsection*{Your turn}
\begin{itemize}
\raggedright
  \item \emph{Say it:} \emph{orior, orīrī}
  \item \emph{Say it:} \emph{orior, orīrī}, once more
  \item \emph{Say it:} a sentence of your own with \emph{orior, orīrī}
  \item \emph{Recall:} say the Latin for to happen, then the Latin for to become, then say \emph{orior, orīrī} again
\end{itemize}

\begin{tcolorbox}[breakable,colback=teal!4,colframe=teal!35!black,title={Before you move on}]
What does orior, orīrī mean? (\textquotedblleft{}To rise; to arise, to begin\textquotedblright{}.) Say it once more.
\end{tcolorbox}
\section[interveniō, intervenīre]{interveniō, intervenīre — to come along, to intervene}

\label{lesson:LA-C298-intervenio}



\begin{quote}
Before the new one: say the Latin for to become, then the Latin for to rise.
\end{quote}

\subsection*{You'll want to know: interveniō, intervenīre}
\textbf{interveniō, intervenīre} — \textquotedblleft{}to come along, to intervene\textquotedblright{}.

\begin{quote}
Pater intervēnit et fābulam interrūpit.
\end{quote}

\begin{quote}
\emph{Father came along and broke off the story.}
\end{quote}

\emph{interveniō, intervenīre, intervēnī, interventum}, often with the dative: \emph{sermōnī intervēnit}, he broke in on the talk.

\begin{cousinweb}
\textbf{interveniō} is built on \emph{veniō}, \textquotedblleft{}to come\textquotedblright{}, with \emph{inter}, \textquotedblleft{}between\textquotedblright{} (Lewis \& Short): to come in during something. English \emph{intervene} comes from it.
\end{cousinweb}

\begin{grammarlens}[title={What you've built}]
One more word for telling, explaining and describing at B1.
\end{grammarlens}

\subsection*{Your turn}
\begin{itemize}
\raggedright
  \item \emph{Say it:} \emph{interveniō, intervenīre}
  \item \emph{Say it:} \emph{interveniō, intervenīre}, once more
  \item \emph{Say it:} a sentence of your own with \emph{interveniō, intervenīre}
  \item \emph{Recall:} say the Latin for to become, then the Latin for to rise, then say \emph{interveniō, intervenīre} again
\end{itemize}

\begin{tcolorbox}[breakable,colback=teal!4,colframe=teal!35!black,title={Before you move on}]
What does interveniō, intervenīre mean? (\textquotedblleft{}To come along, to intervene\textquotedblright{}.) Say it once more.
\end{tcolorbox}
\section[posterus, postera, posterum]{posterus, postera, posterum — next, following}

\label{lesson:LA-C298-posterus}



\begin{quote}
Before the new one: say the Latin for to rise, then the Latin for to come along.
\end{quote}

\subsection*{You'll want to know: posterus, postera, posterum}
\textbf{posterus, postera, posterum} — \textquotedblleft{}next, following\textquotedblright{}.

\begin{quote}
Posterō diē nūntius ad urbem rediit.
\end{quote}

\begin{quote}
\emph{On the next day the messenger returned to the city.}
\end{quote}

\emph{Posterō diē}: the next day, a set phrase in Caesar. The plural noun \emph{posterī} means descendants.

\begin{cousinweb}
\textbf{posterus} is built on \emph{post}, \textquotedblleft{}after, behind\textquotedblright{} (EtymDB): coming after. English \emph{posterity} comes from it.
\end{cousinweb}

\begin{grammarlens}[title={What you've built}]
One more word for telling, explaining and describing at B1.
\end{grammarlens}

\subsection*{Your turn}
\begin{itemize}
\raggedright
  \item \emph{Say it:} \emph{posterus, postera, posterum}
  \item \emph{Say it:} \emph{posterus, postera, posterum}, once more
  \item \emph{Say it:} a sentence of your own with \emph{posterus, postera, posterum}
  \item \emph{Recall:} say the Latin for to rise, then the Latin for to come along, then say \emph{posterus, postera, posterum} again
\end{itemize}

\begin{tcolorbox}[breakable,colback=teal!4,colframe=teal!35!black,title={Before you move on}]
What does posterus, postera, posterum mean? (\textquotedblleft{}Next, following\textquotedblright{}.) Say it once more.
\end{tcolorbox}
\section[(dialogue)]{Words from chapters 291–298 — a review}

\label{lesson:LA-R298-review-291-298}



\begin{quote}
Say the Latin for at first and then.
\end{quote}

\subsection*{Your turn}
Say these aloud:

\begin{itemize}
\raggedright
  \item at first, then, before, not yet, at last
  \item suddenly, at once, straightaway, in an instant, immediately
  \item little by little, gradually, step by step, for a little while, for some time
  \item again and again, immediately after, frequently, mostly, as often as
  \item by night, by day, in the evening, early, a little while ago
  \item a beginning, a start, a way out, a gap, a moment
  \item to begin, I began, to finish, to complete, to break off
  \item to happen, to become, to rise, to come along, next
\end{itemize}

\begin{tcolorbox}[breakable,colback=teal!4,colframe=teal!35!black,title={Before you move on}]
At first? (\textbf{prīmō}.) Then? (\textbf{tunc}.) Before? (\textbf{antequam}.) Not yet? (\textbf{nōndum}.) At last? (\textbf{dēmum}.) Suddenly? (\textbf{repente}.) At once? (\textbf{cōnfestim}.) Straightaway? (\textbf{continuō}.) In an instant? (\textbf{extemplō}.) Immediately? (\textbf{prōtinus}.) Little by little? (\textbf{paulātim}.) Gradually? (\textbf{sēnsim}.) Step by step? (\textbf{gradātim}.) For a little while? (\textbf{paulisper}.) For some time? (\textbf{aliquamdiū}.) Again and again? (\textbf{identidem}.) Immediately after? (\textbf{subinde}.) Frequently? (\textbf{crēbrō}.) Mostly? (\textbf{plērumque}.) As often as? (\textbf{quotiēns}.) By night? (\textbf{noctū}.) By day? (\textbf{interdiū}.) In the evening? (\textbf{vesperī}.) Early? (\textbf{mātūrē}.) A little while ago? (\textbf{dūdum}.) A beginning? (\textbf{initium, initiī}.) A start? (\textbf{prīncipium, prīncipiī}.) A way out? (\textbf{exitus, exitūs}.) A gap? (\textbf{intervallum, intervallī}.) A moment? (\textbf{mōmentum, mōmentī}.) To begin? (\textbf{incipiō, incipere}.) I began? (\textbf{coepī, coepisse}.) To finish? (\textbf{fīniō, fīnīre}.) To complete? (\textbf{perficiō, perficere}.) To break off? (\textbf{intermittō, intermittere}.) To happen? (\textbf{ēveniō, ēvenīre}.) To become? (\textbf{fīō, fierī}.) To rise? (\textbf{orior, orīrī}.) To come along? (\textbf{interveniō, intervenīre}.) Next? (\textbf{posterus, postera, posterum}.)
\end{tcolorbox}
